\chapter{A Concrete Naming Certificate for $\FiniteTailSandwichUp$}
\label{ch:concrete-instances-finite-tail-sandwich-namecert}
\origin{human}

Following Bishop and Bridges, the $\FiniteTailSandwichUp$ packet records the finite tail-squeeze data by which a middle regular rational name is read between lower and upper tail rows at a common dyadic budget. It sits over the arithmetic core of \autoref{ch:concrete-instances-dyadicratcore-namecert}, the finite observation windows of \autoref{ch:concrete-instances-streamname-namecert}, the regular rational sequence interface of \autoref{ch:concrete-instances-regseqrat-namecert}, the terminal real seal of \autoref{ch:concrete-instances-real-namecert}, and the two-sided regular-Cauchy bracketing route of \autoref{ch:concrete-instances-regular-cauchy-sandwich-namecert}. The packet gives Real-and-Completion arguments a displayed finite tail sandwich without using quotient sequence equality, selected subsequences, host limits, or ambient completeness.

\begin{definition}[Finite tail sandwich carrier]
\label{def:finite-tail-sandwich-carrier}
A $\FiniteTailSandwichUp$ carrier is a finite $\BHist$ packet
$$
F=(L,M,U,D,W,R,E,H,C,P,N)
$$
with the following displayed rows:
\begin{enumerate}
\item $L$ is the lower tail row, naming the inspected lower $\RegSeqRatUp$ values and their finite $\StreamNameUp$ windows.
\item $M$ is the middle tail row, naming the inspected middle $\RegSeqRatUp$ values that are to be squeezed.
\item $U$ is the upper tail row, naming the inspected upper $\RegSeqRatUp$ values and their finite $\StreamNameUp$ windows.
\item $D$ is the common $\DyadicRatCoreUp$ budget ledger, recording the shared tail index, tolerance cell, lower--middle comparison, middle--upper comparison, and upper--lower width bound for each displayed read.
\item $W$ is the $\StreamNameUp$ window ledger proving that $L,M,U$ are read on the same finite tail schedule.
\item $R$ is the $\RegSeqRatUp$ readback row through which downstream regular-Cauchy consumers reread the middle row only after the two-sided tail comparisons have been exposed.
\item $E$ is the terminal $\RealUp$ seal row reached after the lower, middle, upper, budget, window, and readback rows have been replayed.
\item $H$ records componentwise $\hsame$ transport for $L,M,U,D,W,R,E$ and the structural rows.
\item $C$ records displayed $\Cont$ replay routes for rereading the same finite sandwich data.
\item $P$ is the $\Pkg$ provenance over $\FiniteTailSandwichUp$, $\DyadicRatCoreUp$, $\StreamNameUp$, $\RegSeqRatUp$, $\RealUp$, $\BHist$, $\hsame$, $\Cont$, and $\NameCert$.
\item $N$ is the local $\NameCert_{\FiniteTailSandwichUp}$ row.
\end{enumerate}
The classifier compares accepted carriers only through these displayed rows. It has no coordinate for host real equality, quotient stream equality, a selected cofinal subsequence, trichotomy, a least upper bound, or ambient $\RealUp$ completeness.
\end{definition}

\begin{theorem}[Finite tail sandwich NameCert obligations]
\label{thm:finite-tail-sandwich-namecert-obligations}
For every accepted $\FiniteTailSandwichUp$ carrier $F=(L,M,U,D,W,R,E,H,C,P,N)$, the local $\NameCert_{\FiniteTailSandwichUp}$ surface consists of carrier admission for the lower, middle, and upper tail rows; common-window exposure through $W$; common dyadic-budget visibility through $D$; regular rational readback through $R$; terminal $\RealUp$ sealing through $E$; componentwise $\hsame$ transport; displayed $\Cont$ replay; $\Pkg$ provenance; and local naming. No obligation may replace these rows by a host limit, quotient sequence equality, selected subsequence, trichotomy, or ambient completeness.
\end{theorem}
\begin{proof}
Project the carrier of \autoref{def:finite-tail-sandwich-carrier}. The analytic rows are exactly $L,M,U,D,W,R,E$, so every admissible obligation must be read from the displayed lower, middle, upper, budget, window, readback, and seal coordinates.

The common-tail condition is enforced by $W$ and $D$. The window row fixes the finite tail schedule shared by the three regular rational rows, while the dyadic ledger records the lower--middle and middle--upper comparisons at the same tolerance cell.

The structural rows $H,C,P,N$ preserve this inventory. Transport keeps the same component rows visible, continuation replay rereads them in the displayed order, provenance names the packet source, and the local name records no additional analytic coordinate. Hence the listed rows exhaust the local NameCert surface and exclude the non-constructive escape routes named in the statement.
\end{proof}

\begin{theorem}[Finite tail sandwich RegSeqRat handoff]
\label{thm:finite-tail-sandwich-regseqrat-handoff}
Every $\RegSeqRatUp$ or $\RealUp$ consumer of the middle row in an accepted $\FiniteTailSandwichUp$ packet factors through the displayed route
$$
L\longmapsto M\longmapsto U\longmapsto D\longmapsto W\longmapsto R\longmapsto E\longmapsto H\longmapsto C\longmapsto P\longmapsto N .
$$
In particular, the handoff exposes a finite lower--middle--upper tail squeeze at a common dyadic budget, not a host limit, quotient sequence equality, selected subsequence, or ambient completeness principle.
\end{theorem}
\begin{proof}
Begin with the lower, middle, and upper rows $L,M,U$. They supply only finite $\RegSeqRatUp$ reads and finite $\StreamNameUp$ windows, so the middle row is never detached from its two bounding rows.

Next read $D$ and $W$. The dyadic ledger stores the two-sided comparison and width budget at the inspected tolerance, while the window ledger records that all three rows are read on the same finite tail schedule.

Only after those rows are visible can a consumer reach $R$ and $E$. The readback row makes the middle regular rational name available to downstream regular-Cauchy consumers, and the terminal seal exposes the Real-facing boundary. The structural rows $H,C,P,N$ then transport, replay, source, and name precisely the same route, so no host limit or completeness coordinate is introduced.
\end{proof}

The chapter is the finite-tail sibling of the broader regular-Cauchy sandwich certificate in \autoref{ch:concrete-instances-regular-cauchy-sandwich-namecert}. That sibling records the two-sided regular-Cauchy bracketing route as a full certificate; the present packet isolates the common-tail budget, window, and middle-row readback used when a later Real-and-Completion argument only needs the finite squeeze data.

\closureat{FiniteTailSandwichUp}{seedStr}

\begin{closurestatus}{\FiniteTailSandwichUp}
  \constructivestory{A $\FiniteTailSandwichUp$ packet is generated from lower, middle, and upper RegSeqRat tail rows, a common DyadicRatCore budget ledger, StreamName windows, RegSeqRat readback, a terminal Real seal, hsame transport, Cont replay, Pkg provenance, and a local NameCert row.}
  \theoryclosure{\seedClosure}
  \scopeclosed{\autoref{def:finite-tail-sandwich-carrier}, \autoref{thm:finite-tail-sandwich-namecert-obligations}, and \autoref{thm:finite-tail-sandwich-regseqrat-handoff} seed the finite lower--middle--upper tail squeeze surface for Real-and-Completion consumers.}
  \formalstatus{\unformalizedV}
  \bridgestatus{none}
  \notclaimed{This chapter does not close Lean verification, bridge checking, quotient sequence equality, selected subsequences, host limits, trichotomy, least-upper-bound reasoning, or ambient $\RealUp$ completeness.}
  \upgradepath{Obligation closure requires checked carrier admission, common-tail window exposure, dyadic two-sided budget visibility, RegSeqRat readback, Real seal routing, transport, replay, provenance, and local naming for BEDC.Derived.FiniteTailSandwichUp.}
\end{closurestatus}
